% !TEX root = ../main.tex
%-----------------------------------------------------------------------
%  PART II -- the neural-network surrogate design.
%  \input by ../main.tex -- not compiled on its own.
%-----------------------------------------------------------------------
\part{A neural-network surrogate for the sensitivity fields}

%-----------------------------------------------------------------------
\section{Summary at a glance}
%-----------------------------------------------------------------------

\begin{center}
\begin{tabularx}{\textwidth}{@{}lL@{}}
\toprule
\textbf{Target} & Emulate the map
$\big(\bar{\mathbf{x}},\,\boldsymbol{\theta},\,\tau,\,c\big)\mapsto
\partial J/\partial\boldsymbol{\theta}$, i.e.\ the adjoint operator itself,
not the forward model.\\
\midrule
\textbf{Primary output} & $\partial J/\partial\kv(x,y,z)$ on the DINO grid --
one $51\times198\times36$ field, $\sim\!3.3\times10^{5}$ wet points.\\
\textbf{Primary inputs} & Forward state (T, S, U, V, W, $\rho'$), derived
stratification diagnostics, the parameter fields themselves, grid geometry and
masks, a cost-section channel, and lead time $\tau$.\\
\textbf{Loss} & Volume-weighted Huber on an $\operatorname{asinh}$-transformed,
per-lead-time-normalised field $+$ a separate log-amplitude head $+$ a
pattern-cosine term $+$ a finite-difference directional-derivative term.\\
\textbf{Architecture} & 3-D U-Net with FiLM conditioning on $(\tau,
\boldsymbol{\theta}_{\text{scalar}})$; POD--ridge reduced-basis model as the
baseline it has to beat; Fourier neural operator as the first alternative.\\
\textbf{Data} & 200--400 member Sobol ensemble over 6--8 parameters
$\times$ 8--12 background states $\times$ 3--5 cost sections. What we hold now
is enough for a baseline and for building the pipeline, but not to train a
parametric surrogate.\\
\textbf{Biggest lever on cost} & Turning off \code{useGrdchk} for ensemble
members. A 30-day adjoint currently executes $18{,}622$ forward steps for the
$1{,}440$ it integrates ($12.9\times$).\\
\textbf{Biggest open question} & How far back in lead time the labels stay
meaningful. I do not yet know where that horizon sits for DINO, and it sets
the ensemble cost.\\
\bottomrule
\end{tabularx}
\end{center}

%-----------------------------------------------------------------------
\section{What the adjoint output already gives us}
\label{sec:adjoutput}
%-----------------------------------------------------------------------

I have gone through the output we have and measured it rather than assuming,
because the design that follows depends on these numbers. The grid and cost
functional are as given in Section~\ref{sec:modelconfig}.

\paragraph{Adjoint output.}
The reverse sweep dumps eleven adjoint fields at \code{adjDumpFreq}. For the
reference five-year run (job 28486, from the 180-year pickup, \code{visc2x})
that is $366$ dumps at five-day spacing:

\begin{center}
\begin{tabularx}{\textwidth}{@{}Y{0.7}Y{1.1}Y{1.2}@{}}
\toprule
\textbf{Class} & \textbf{Fields} & \textbf{Note}\\
\midrule
State / initial condition & \code{ADJtheta}, \code{ADJsalt},
  \code{ADJuvel}, \code{ADJvvel}, (\code{ADJwvel})
  & 3-D, \code{float32}, $51\times198\times36$. \code{ADJwvel} is dumped but
    \code{xx\_wvel} is not among the controls declared in \code{data.ctrl} ---
    $w$ is diagnosed from continuity, not an initial condition --- so it is not
    a prediction target and is bracketed here\\
Surface forcing & \code{ADJtaux}, \code{ADJtauy}, \code{ADJqnet},
  \code{ADJempmr}, \code{ADJqsw}
  & 2-D, \code{float32}\\
\textbf{Parametric} & \code{ADJdiffkr}
  & 3-D --- \emph{the target}: $\partial J/\partial\kv$\\
Parametric, not active & \code{ADJbtdrg}, \code{ADJkapgm}, \code{ADJkapredi}
  & dumped by \code{addummy\_in\_stepping.F} but need controls declared and,
    for the last two, GM/Redi enabled\\
\midrule
Accumulated controls & \code{adxx\_*} & \code{float64}, written once per run\\
\bottomrule
\end{tabularx}
\end{center}

\paragraph{Two structural facts I checked numerically.}
Both of these affect how the problem should be posed.

\begin{enumerate}
  \item \code{ADJdiffkr.<iter>} is the gradient \emph{accumulated so far} over
        the reverse sweep, not an independent snapshot at each lag. Its RMS
        over wet points rises monotonically as the sweep runs backwards, from
        $6.55\times10^{-4}$ at the end of the forward window ($\tau\approx0$)
        to $4.30\times10^{-1}$ at the start ($\tau=5\,\mathrm{yr}$), a factor
        of $656$ or about $2.8$ decades.
  \item The earliest-iteration dump is the total gradient. Compared against
        \code{adxx\_diffkr} the correlation is $1.000000$, the RMS agrees to
        five digits, and the median pointwise ratio is $1.0000$. So
        \code{adxx\_diffkr} is $\partial J/\partial\kv$, and the dump sequence
        is its accumulation curve in lead time.
\end{enumerate}

\paragraph{The three classes behave differently in lead time.}
Only \code{ADJdiffkr} accumulates. On the same run, running backwards over the
five years, \code{ADJdiffkr} RMS grows $656\times$ while \code{ADJtheta} decays
by about $5\times$ and \code{ADJtaux} by about $40\times$. The latter two are
adjoint states being propagated and damped, not integrals. The three also sit
about $10^{5}$ apart in absolute magnitude ($\sim\!10^{-1}$, $\sim\!10^{-4}$
and $\sim\!10^{-6}$), which is why each output head needs its own
normalisation (Section~\ref{sec:norm}).

\paragraph{Dynamic range and sign structure.}
Over wet points, $|\partial J/\partial\kv|$ covers roughly fifteen decades: the
$1^{\text{st}}$, $50^{\text{th}}$ and $99^{\text{th}}$ percentiles of
$\log_{10}|g|$ are $-14.7$, $-3.0$ and $-0.04$. The sign is close to balanced,
$48.6\,\%$ negative against $48.5\,\%$ positive. This one measurement is what
drives the normalisation and loss design in Sections~\ref{sec:norm}
and~\ref{sec:loss}. An unweighted mean-squared error on the raw field fits the
few hundred cells nearest the section and ignores the rest of the basin.

\paragraph{Forward state.}
The \code{data.diagnostics} streams give us, at $30.5$-day spacing ($60$
snapshots over five years): \code{dynDiag} (THETA, SALT, PHIHYD, RHOAnoma,
UVEL, VVEL, WVEL), \code{surfDiag} (ETAN, TRELAX, SRELAX, MXLDEPTH, PHIBOT),
\code{atmDiag} (oceTAUX, oceTAUY, oceQnet) and \code{viscDiag} (VISCAHZ,
VISCAHD). Grid and mask files (\code{hFacC}, \code{Depth}, \code{XC},
\code{YC}, \code{RC}, \code{DRF}, \code{DXC}, \dots) are written into every run
directory.

\paragraph{Existing runs.}
There are fifteen DINO adjoint run directories on scratch: four five-year MPI
runs, one ten-year, one thirty-day, five at 180 or 360 days, and four serial
runs between ten days and two years. Between them they cover three background
states (from rest, and from the 50- and 180-year pickups) and four viscosity
treatments. Two of the five-year runs used \code{adjViscBoost}, which means
they came from a different adjoint operator than the others.

None of this was generated as a designed ensemble, and along the parameters we
care about the runs vary one axis at a time rather than jointly, so for
learning a parametric map the effective sample size is three or four. That is
enough to build and debug the ingest pipeline, characterise the target
statistics and fit the reduced-basis baseline. It is not enough to train the
surrogate. Generating the ensemble is the critical path here, and everything in
Phase~1 exists to make it affordable.

%-----------------------------------------------------------------------
\section{What ``parametric sensitivity'' means here}
%-----------------------------------------------------------------------

The control vector in \code{input\_tap/data.ctrl} holds three kinds of thing,
and only one of them is parametric:

\begin{center}
\begin{tabularx}{\textwidth}{@{}Y{0.55}Y{1.25}Y{1.2}@{}}
\toprule
\textbf{Class} & \textbf{Controls} & \textbf{Status}\\
\midrule
Initial state & \code{xx\_theta}, \code{xx\_salt}, \code{xx\_uvel},
  \code{xx\_vvel} & active, 3-D, time-invariant\\
Surface forcing & \code{xx\_tauu}, \code{xx\_tauv}, \code{xx\_qnet},
  \code{xx\_empmr}, \code{xx\_qsw}, \code{xx\_uwind}, \code{xx\_vwind},
  \code{xx\_fu}, \code{xx\_fv} & active, 2-D, time-varying\\
\textbf{Parametric} & \code{xx\_diffkr} & \textbf{active} --- the only
  parametric control in the vector at present\\
\bottomrule
\end{tabularx}
\end{center}

The parameters I would like to add, and how they currently reach the model:

\begin{center}
\begin{tabularx}{\textwidth}{@{}Y{0.7}Y{0.75}Y{1.55}@{}}
\toprule
\textbf{Parameter} & \textbf{Enters via} & \textbf{To make it a control}\\
\midrule
Vertical diffusivity $\kv(x,y,z)$ & \code{diffKrFile} (PARM05)
  & already done (\code{xx\_diffkr})\\
Lateral viscosity $\Ah(x,y)$ & \code{viscAhDfile}, \code{viscAhZfile}
  & add a \code{genarr2d} control; it is already a read-in field\\
Grid viscosity \code{viscAhGrid} & PARM01 scalar & scalar: finite differencing is
  competitive\\
Vertical viscosity \code{viscAr} & PARM01 scalar & scalar\\
Bottom drag \code{bottomDragQuadratic} & PARM01 scalar
  & \code{ADJbtdrg} is already dumped; needs \code{xx\_bottomdrag}\\
Convective adjustment \code{ivdc\_kappa} & PARM01 scalar & scalar\\
Lateral T/S diffusivity \code{diffKhT/S} & PARM01 scalar & scalar\\
GM / Redi $\kappa_{\mathrm{GM}},\kappa_{\mathrm{Redi}}$
  & \code{data.gmredi} & \code{useGMRedi=.FALSE.} at present; enabling it
    changes the physics and needs a fresh spin-up\\
\bottomrule
\end{tabularx}
\end{center}

\paragraph{Where the surrogate earns its keep.}
For a scalar parameter the adjoint returns one number that a pair of forward
runs would also give, so a surrogate adds little. The case for it rests on the
field-valued parameters: $\kv(x,y,z)$ with $3.3\times10^{5}$ wet components and
$\Ah(x,y)$ with $9.6\times10^{3}$, where finite differencing is not an option
and the adjoint is the only practical route. So $\kv$ is the primary target
here, $\Ah$ the first extension, and the scalars come in as conditioning inputs
rather than as things we try to predict.

%-----------------------------------------------------------------------
\section{The three sensitivity classes, and why one network covers all of them}
\label{sec:classes}
%-----------------------------------------------------------------------

One reverse sweep writes all three classes of sensitivity. That is the piece of
economy this project rests on: a network trained on the parametric target can
carry initial-state and surface-forcing heads for almost no extra data, since
the labels are already sitting on disk. What differs between the classes is the
shape of the target and what the lead-time axis means.

\begin{center}
\begin{tabularx}{\textwidth}{@{}Y{0.60}Y{0.75}Y{0.70}Y{1.95}@{}}
\toprule
\textbf{Class} & \textbf{Adjoint fields} & \textbf{Target shape} &
\textbf{What the $\tau$ axis means, and how it is emulated}\\
\midrule
\textbf{Parametric}
  & \code{ADJdiffkr}; total in \code{adxx\_diffkr}
  & $51{\times}198{\times}36$
  & \emph{Accumulation.} The gradient integrated over
    $[t_{\text{end}}-\tau,\,t_{\text{end}}]$, so it is monotone in $\tau$ and
    the $\tau_{\max}$ limit is what we are after. A 3-D decoder head, with
    monotonicity available as a constraint ($\mathcal{L}_{\text{acc}}$).\\
\textbf{Initial state}
  & \code{ADJtheta}, \code{ADJsalt}, \code{ADJuvel}, \code{ADJvvel}
  & $4{\times}51{\times}198{\times}36$
  & \emph{Propagation.} The adjoint state at $t_{\text{end}}-\tau$: sensitivity
    to a perturbation applied \emph{at that time}, not an integral. Same tensor
    shape as the parametric target, so it reuses that decoder with more output
    channels. The cheapest extension available.\\
\textbf{Surface forcing}
  & \code{ADJtaux}, \code{ADJtauy}, \code{ADJqnet}, \code{ADJempmr},
    \code{ADJqsw}
  & $5{\times}51{\times}198$ per $\tau$
  & \emph{Time-resolved.} Sensitivity to forcing applied at
    $t_{\text{end}}-\tau$; the physical object is the space--time stack over
    $\tau$. Since $\tau$ is a conditioning input already, one 2-D head predicts
    the whole stack. The cheapest head by some distance.\\
\bottomrule
\end{tabularx}
\end{center}

\paragraph{Each head needs its own normalisation.}
The classes are about five orders apart in magnitude ($\sim\!10^{-1}$,
$\sim\!10^{-4}$, $\sim\!10^{-6}$), and their lead-time trends run in opposite
directions: the parametric target grows because it accumulates, while the other
two decay as they propagate. Sharing one loss scale would let the parametric
head dominate, and would fit a single amplitude trend to three different
functions. So Section~\ref{sec:norm} applies per head, and the loss in
Section~\ref{sec:loss} is summed over heads with per-class weights.

\paragraph{A configuration point about the forcing controls.}
\code{adxx\_tauu} is written with \code{nrecords = 1}. No
\code{xx\_gentim2d\_period} is set in \code{input\_tap/data.ctrl}, so although
the forcing controls are declared time-varying (\code{xx\_gentim2d}), the
accumulated gradient collapses to a single 2-D map instead of a time series.
Time-resolved forcing sensitivity is available today only through the $366$
\code{ADJtaux} dumps, which is fine for training. But \code{adxx\_tauu} should
not be read as a time series, and setting a control period is the change to
make if we want the time-resolved accumulated gradient.

%-----------------------------------------------------------------------
\section{Three ways to pose the surrogate}
%-----------------------------------------------------------------------

Let $\bar{\mathbf{x}}$ be the forward background state, $\boldsymbol{\theta}$
the parameter set, $\tau$ the lead time (time before the cost evaluation), and
$c$ the cost-functional specification. Write the target as
$g_\tau = \partial J/\partial\boldsymbol{\theta}\big|_{[t_{\text{end}}-\tau,\,t_{\text{end}}]}$.

\paragraph{F1 --- Direct operator emulation (what I would start with).}
Learn $G_\phi:(\bar{\mathbf{x}},\boldsymbol{\theta},\tau,c)\mapsto\hat{g}_\tau$
directly. The labels are exactly what the adjoint already writes, and nothing
gets differentiated, so there is no compounding of a fitted function's
derivative error. This is the version the existing data support with no change
to the pipeline.

\paragraph{F2 --- Emulate the scalar cost, then differentiate the network.}
Learn $J_\phi(\boldsymbol{\theta};\bar{\mathbf{x}})$ from forward runs alone and
get $\partial J/\partial\boldsymbol{\theta}$ by differentiating the network.
The training data would be cheap, since no adjoint runs are needed at all. The
problem is the familiar one: a network fitted to function values can be
accurate in $J$ and badly wrong in $\nabla J$, especially in a
high-dimensional parameter space where $J$ is nearly flat in most directions.
I would not do this on its own.

\paragraph{F3 --- Sobolev / gradient-matched training (the natural extension).}
Learn $J_\phi$ with a loss on both the value and the gradient,
\begin{equation}
\mathcal{L} \;=\;
\big\|J_\phi - J\big\|^{2}
\;+\;
\lambda\,\Big\|\nabla_{\boldsymbol{\theta}} J_\phi
- \underbrace{\partial J/\partial\boldsymbol{\theta}}_{\text{adjoint label}}\Big\|^{2},
\end{equation}
using the adjoint output to supervise the derivative. That gives one object
consistent in both value and gradient, which can be dropped straight into an
optimiser. It needs the same ensemble as F1, so it follows F1 rather than
competing with it.

\paragraph{A structural point worth using.}
The adjoint field solves the \emph{linear} adjoint equation: it is linear in
the terminal seed (the derivative of the cost functional with respect to the
model state at $t_{\text{end}}$) and nonlinear in
$(\bar{\mathbf{x}},\boldsymbol{\theta})$. A surrogate written as
\begin{equation}
\hat{g}_\tau \;=\; \mathcal{A}_\phi\!\left(\bar{\mathbf{x}},
\boldsymbol{\theta},\tau\right)\big[\,s_c\,\big],
\qquad s_c=\partial J/\partial\mathbf{x}\big|_{t_{\text{end}}},
\end{equation}
with the seed $s_c$ entering as an input field rather than being absorbed into
the weights, respects that linearity and generalises across cost functionals
without being asked to. Given that the section indices are compiled in, this is
how the surrogate could answer questions the compiled model cannot answer
without a rebuild. In practice it means feeding $s_c$ as an input channel and
putting several sections in the training ensemble.

%-----------------------------------------------------------------------
\section{Inputs}
\label{sec:inputs}
%-----------------------------------------------------------------------

All fields sit on the native $51\times198\times36$ grid ($51\times198$ for the
2-D ones). Channels marked \emph{derived} are computed at ingest.

\begin{center}
\begin{tabularx}{\textwidth}{@{}Y{0.6}Y{1.15}cY{1.25}@{}}
\toprule
\textbf{Group} & \textbf{Channels} & \textbf{$n$} & \textbf{Source / rationale}\\
\midrule
3-D state & THETA, SALT, UVEL, VVEL, WVEL, RHOAnoma & 6
  & \code{dynDiag}, interpolated to the lead time\\
3-D derived & $N^{2}$, $\partial\bar{T}/\partial z$, vertical shear
  $|\partial\mathbf{u}/\partial z|$, $|\nabla_h b|$ & 4
  & \emph{derived.} $\kv$ acts on the vertical tracer gradient, so
    $\partial\bar{T}/\partial z$ and $N^{2}$ should be the leading local
    predictors of $\partial J/\partial\kv$. Supplying them is cheaper than
    making the network rediscover a finite difference\\
2-D state & ETAN, MXLDEPTH, oceTAUX, oceTAUY, oceQnet & 5
  & \code{surfDiag}, \code{atmDiag}; broadcast, or through a separate 2-D
    encoder\\
Parameter fields & $\kv(x,y,z)$, $\Ah(x,y)$, VISCAHD, VISCAHZ & 4
  & the point of the exercise; $\Ah$ broadcast in the vertical\\
Static geometry & \code{hFacC}, \code{Depth}, \code{maskC}, $y$ (latitude),
  $z$ (\code{RC}), \code{DRF}, cell area & 7
  & the masks do real work here: $8.6\,\%$ of cells are land and have to be
    excluded everywhere, not zeroed\\
Cost specification & section indicator; terminal seed $s_c$ & 2
  & what makes the surrogate re-targetable across cost functionals\\
\bottomrule
\end{tabularx}
\end{center}

\paragraph{Scalar conditioning.}
Nine further quantities go in through FiLM layers rather than as broadcast
channels, since they carry no spatial structure: the lead time $\tau$ (as
Fourier features), \code{viscAr}, \code{bottomDragQuadratic},
\code{ivdc\_kappa}, \code{diffKhT/S}, \code{viscAhGrid},
\textbf{\code{viscFacInAd}}, the seasonal phase of $t_{\text{end}}$, and
$\text{nIter}_0$ as a proxy for spin-up age. In total the network sees about
$28$ field channels and $9$ scalars.

\paragraph{\code{viscFacInAd} has to be an input, not a constant.}
The \code{adjViscBoost} configuration runs the adjoint sweep with inflated
viscosity and diffusivity relative to the forward model
(\code{viscFacInAd}${}=10$ against \code{viscFacInFw}${}=1$), which is what
keeps a long adjoint from blowing up. It also changes the operator being
emulated. Mixing boosted and unboosted sweeps in the training set without
telling the network which is which means training on two different labelling
functions. Either hold it fixed across the ensemble or, better, treat it as a
conditioning scalar and sample it, which has the side benefit of letting us
measure how much the boost distorts the sensitivity.

%-----------------------------------------------------------------------
\section{Outputs}
%-----------------------------------------------------------------------

\begin{center}
\begin{tabularx}{\textwidth}{@{}Y{0.65}Y{0.7}Y{1.65}@{}}
\toprule
\textbf{Head} & \textbf{Shape} & \textbf{Content}\\
\midrule
Primary (pattern) & $51\times198\times36$
  & normalised, $\operatorname{asinh}$-transformed
    $\partial J/\partial\kv$ at lead time $\tau$\\
Primary (amplitude) & scalar
  & $\log_{10}$ of the volume-weighted RMS of the field at that lead time\\
Auxiliary 3-D & $4\times51\times198\times36$
  & \code{ADJtheta}, \code{ADJsalt}, \code{ADJuvel}, \code{ADJvvel} --- labels
    we get for free, and training on them regularises the shared encoder\\
Auxiliary 2-D & $5\times51\times198$
  & \code{ADJtaux}, \code{ADJtauy}, \code{ADJqnet}, \code{ADJempmr},
    \code{ADJqsw}\\
Total gradient & $51\times198\times36$
  & \code{adxx\_diffkr}, the $\tau=\tau_{\max}$ limit --- the quantity an
    optimiser actually consumes\\
\bottomrule
\end{tabularx}
\end{center}

Since the dump sequence is an accumulation
(Section~\ref{sec:adjoutput}), the total-gradient head
and the sequence of pattern heads are not independent, and the consistency
between them is a training signal we get for nothing (the
$\mathcal{L}_{\text{acc}}$ term below). The auxiliary heads are auxiliary only
in the sense that the parametric target drives the design. They are full
sensitivity products in their own right, and Section~\ref{sec:classes} says
what each of them means.

%-----------------------------------------------------------------------
\section{Normalisation}
\label{sec:norm}
%-----------------------------------------------------------------------

The target field is sign-symmetric, spans fifteen decades, and shifts in
overall amplitude by nearly three decades across the lead-time range. Standard
normalisation handles none of the three, which is why I have given this its own
section. What I would do:

\begin{enumerate}
  \item \textbf{Mask.} Every loss and every statistic restricted to wet points,
        $\code{hFacC}>0$. Land excluded, not set to zero.
  \item \textbf{Volume weights.} Define
        $w_{ijk} = \code{hFacC}_{ijk}\,\code{DRF}_k\,\code{rA}_{ij}$, normalised
        so $\sum w = 1$. The DINO vertical grid stretches from $10\,\mathrm{m}$
        to $600\,\mathrm{m}$, so without volume weighting the deep ocean counts
        for almost nothing and the loss follows the surface layers.
  \item \textbf{Per-(field, lead-time) robust scale.}
        $s_\tau = \big(\sum_{ijk} w_{ijk}\,g_{ijk}^{2}\big)^{1/2}$, computed
        over the training set at each lead time.
  \item \textbf{Signed logarithmic transform.}
        $\tilde{g} = \operatorname{asinh}\!\big(g/(\alpha s_\tau)\big)$ with
        $\alpha\approx10^{-2}$: linear near zero, logarithmic in the tails, odd
        and smooth. A $\operatorname{sign}(g)\log(1+|g|)$ construction has the
        same properties, but $\operatorname{asinh}$ has the cleaner derivative
        for backpropagation.
  \item \textbf{Amplitude/pattern split.} The network predicts the normalised
        pattern and $\log_{10}s_\tau$ separately. That takes the lead-time
        amplitude trend --- a smooth monotone function of $\tau$, different for
        each sensitivity class --- out of the spatial problem, which is the hard
        part. Reconstruct as
        $\hat{g} = \alpha\,10^{\widehat{\log_{10}s_\tau}}\sinh(\hat{\tilde{g}})$.
  \item \textbf{Winsorise} at the $99.9^{\text{th}}$ weighted percentile before
        training, and report metrics both with and without the clipping.
\end{enumerate}

Input channels get standardised per channel using training-set statistics over
wet points, except the parameter channels, which go to $\log_{10}$ first
($\kv$ runs from $10^{-6}$ to $5\times10^{-4}$ $\mathrm{m^{2}\,s^{-1}}$ under
its \code{data.ctrl} bounds).

%-----------------------------------------------------------------------
\section{Loss function}
\label{sec:loss}
%-----------------------------------------------------------------------

One composite objective, with each term doing something the others do not:
\begin{equation}
\mathcal{L}
= \lambda_{p}\,\mathcal{L}_{\text{pattern}}
+ \lambda_{m}\,\mathcal{L}_{\text{amp}}
+ \lambda_{c}\,\mathcal{L}_{\cos}
+ \lambda_{a}\,\mathcal{L}_{\text{acc}}
+ \lambda_{d}\,\mathcal{L}_{\text{dir}}
+ \lambda_{x}\,\mathcal{L}_{\text{aux}} .
\end{equation}

\paragraph{$\mathcal{L}_{\text{pattern}}$ --- weighted Huber in transformed space.}
\begin{equation}
\mathcal{L}_{\text{pattern}}
= \sum_{ijk\in\Omega_{\text{wet}}} w_{ijk}\;
  \mathrm{Huber}_{\delta}\!\big(\hat{\tilde{g}}_{ijk}-\tilde{g}_{ijk}\big).
\end{equation}
Huber rather than squared error, because the residuals stay heavy-tailed even
after the $\operatorname{asinh}$ transform; $\delta\approx1$ in transformed
units.

\paragraph{$\mathcal{L}_{\text{amp}}$ --- amplitude head.}
$\big|\widehat{\log_{10}s_\tau}-\log_{10}s_\tau\big|^{2}$. Plain squared error
is fine here, since this quantity varies smoothly over about three decades.

\paragraph{$\mathcal{L}_{\cos}$ --- pattern agreement.}
\begin{equation}
\mathcal{L}_{\cos}
= 1 - \frac{\langle \hat{g},\,g\rangle_{w}}
            {\|\hat{g}\|_{w}\,\|g\|_{w}},
\qquad
\langle a,b\rangle_{w}=\sum_{ijk}w_{ijk}\,a_{ijk}b_{ijk}.
\end{equation}
This is the term that matches what we actually want out of the surrogate: the
shape of the sensitivity map, and for optimisation the direction of the
gradient. A surrogate with $10\,\%$ amplitude error and cosine similarity
$0.99$ is useful to us; one with the errors the other way round is not.

\paragraph{$\mathcal{L}_{\text{acc}}$ --- accumulation consistency.}
Penalise disagreement between the predicted total-gradient head and the
predicted sequence of accumulated fields, and require the predicted amplitude
to be monotone in $\tau$. This encodes a property of the data I measured rather
than assumed.

\paragraph{$\mathcal{L}_{\text{dir}}$ --- finite-difference anchor.}
For a small set of perturbation directions $\delta\boldsymbol{\theta}$ with
finite-difference cost changes measured from paired forward runs,
\begin{equation}
\mathcal{L}_{\text{dir}}
= \Big|\;\langle \hat{g},\,\delta\boldsymbol{\theta}\rangle
  - \frac{J(\boldsymbol{\theta}+\epsilon\,\delta\boldsymbol{\theta})
        - J(\boldsymbol{\theta}-\epsilon\,\delta\boldsymbol{\theta})}
         {2\epsilon}\;\Big|^{2}.
\end{equation}
This is the only term tied to something directly observable rather than to the
adjoint's own output, so it is what stops us inheriting a systematic error in
the adjoint. It is also cheap: a handful of directions per parameter setting,
reusing forward runs the ensemble produces anyway. It is the gradient check in
another form --- see Section~\ref{sec:val}.

\paragraph{$\mathcal{L}_{\text{aux}}$.} The same pattern loss on the auxiliary
state and forcing heads, down-weighted ($\lambda_{x}\approx0.1$), where the
point is regularising the shared encoder.

\paragraph{Starting weights.}
$\lambda_{p}=1$, $\lambda_{m}=1$, $\lambda_{c}=1$, $\lambda_{a}=0.1$,
$\lambda_{d}=0.1$, $\lambda_{x}=0.1$, then tuned on validation. I would expect
$\lambda_{c}$ to want raising once the amplitude head has converged.

%-----------------------------------------------------------------------
\section{Architecture}
%-----------------------------------------------------------------------

\paragraph{The baseline comes first.}
Before any network: a proper-orthogonal-decomposition reduced-basis model.
Project the training sensitivity fields onto their leading $r\approx20$--$100$
volume-weighted POD modes and regress the coefficients on
$(\boldsymbol{\theta},\tau,\text{state summary})$ with ridge regression or
gradient-boosted trees. Adjoint fields from one configuration are strongly
correlated, and a few dozen modes may well capture most of the parametric
variation. If they do, that is the result, and a better one than a network,
because we can look at the modes and say what they are. The network has to beat
this to be worth building.

\paragraph{Primary: conditioned 3-D U-Net.}
Four resolution levels, GroupNorm, SiLU activations, residual blocks; FiLM
conditioning on $(\tau,\boldsymbol{\theta}_{\text{scalar}})$ at every level;
masked convolutions so land never contributes. Three things about this grid
need care:

\begin{itemize}[nosep]
  \item The vertical is $36$ strongly stretched levels, and isotropic
        $3\times3\times3$ kernels quietly assume it is not. Either use
        anisotropic kernels ($3\times3$ horizontal with separate 1-D vertical
        mixing) or fold depth into the channel dimension and work on a
        $51\times198$ field with $36C$ channels, which on ocean grids this size
        is often both cheaper and more accurate.
  \item The basin is closed, pole to pole, and periodic in neither direction.
        Pad zero or replicate at the walls, not circular.
  \item Four downsampling levels on a $198$-point meridional axis give a
        receptive field that spans the basin, which the physics needs:
        sensitivity travels upstream along advective pathways over thousands of
        kilometres, so a shallow local CNN will miss the remote structure
        systematically.
\end{itemize}

\paragraph{Alternative worth testing: Fourier neural operator.}
An FNO in the horizontal with a local vertical stencil represents the nonlocal,
basin-scale structure directly, which is exactly what makes this hard for
convolutions, and it is resolution-flexible if we ever move to a different
grid. Worth one comparison run, but not where I would start.

\paragraph{Lead time: conditioning or rollout.}
Two options. (a) One network with $\tau$ as a conditioning variable, predicting
any lead time in a single shot: simple, no error accumulation, and where I would
begin. (b) A reverse-time propagator
$\hat{g}_{\tau+\Delta}=R_\phi(\hat{g}_\tau,\bar{\mathbf{x}},\boldsymbol{\theta})$
applied autoregressively, which emulates the adjoint propagator itself, is the
most faithful to the mathematics, and can extrapolate past the training
horizon --- but it accumulates error and is harder to train. Phase 4.

\paragraph{Scale.} $5$--$20$M parameters. A $51\times198\times36$ field with
about $30$ channels fits comfortably on one A100/GH200 at moderate batch size.
Mixed precision, but with the loss reduction kept in \code{float32} given the
dynamic range.

%-----------------------------------------------------------------------
\section{Training data: the full parameter ensemble}
\label{sec:trainingdata}
%-----------------------------------------------------------------------

\paragraph{Design.}
A Sobol or Latin-hypercube design over the parameter box, crossed with
background states and cost sections:

\begin{center}
\begin{tabularx}{\textwidth}{@{}Y{0.85}Y{0.55}Y{1.6}@{}}
\toprule
\textbf{Axis} & \textbf{Levels} & \textbf{Range / note}\\
\midrule
$\kv$ background level & continuous & $10^{-6}$ to $5\times10^{-4}$
  $\mathrm{m^{2}\,s^{-1}}$, per the \code{xx\_diffkr} bounds already in
  \code{data.ctrl}\\
$\kv$ profile shape & 2 params & e.g.\ surface value and an
  exponential/tanh depth scale\\
$\Ah$ scaling & continuous & multiples of the reference
  \code{dino\_viscAhD.bin} field ($\mathrm{d}x_C\cdot0.27/2$)\\
\code{viscAr}, \code{bottomDragQuadratic}, \code{ivdc\_kappa},
  \code{diffKhT/S} & 4 scalars & factor $\sim\!3$ either side of current values\\
Background state & 8--12 & pickups spanning years 50--200 of the completed
  200-year spin-up (run 28463), plus seasonal phase\\
Cost section & 3--5 & different latitudes and depth ranges\\
\bottomrule
\end{tabularx}
\end{center}

Target $200$--$400$ members. I would hold out whole corners of the parameter
box, whole background states, and at least one cost section, so that what gets
tested is generalisation rather than interpolation inside a member.

\paragraph{Cost, and what can be done about it.}
A five-year adjoint costs about $622$ core-hours today
(Section~\ref{sec:compute}), so $300$ of them would
be $\sim\!1.9\times10^{5}$ core-hours. That is not a reasonable thing to ask
for. Three reductions, largest first:

\begin{enumerate}
  \item \textbf{Turn off \code{useGrdchk} for ensemble members.} Every adjoint
        job currently runs the finite-difference gradient check as well. On the
        30-day run, \code{FORWARD\_STEP} was entered $18{,}622$ times for a run
        integrating $1{,}440$ timesteps, a factor of $12.9$, most of it the
        check's perturbed forward runs on top of checkpoint recomputation. This
        is the largest saving available and also the easiest.
  \item \textbf{Shorten most members to one or two years.} Sensitivity is more
        trustworthy at shorter lead times in any case
        (Section~\ref{sec:risks}). Keep a few five-year members to test
        lead-time extrapolation.
  \item \textbf{Tapenade \code{-nocheckpoint} tuning.}
        \code{tools/tapenade\_profiling/} already holds the 64-routine list
        from the earlier profiling work, and c69m's \code{genmake2} takes
        \code{-tap\_extra}, so this is a build flag now rather than a patched
        script.
\end{enumerate}

\paragraph{What (2) does not save.}
Shortening the adjoint window does not shorten the forward leg. Any member drawn
at a new parameter setting has to be re-equilibrated before its adjoint window
opens, and at the measured forward rate ($0.0332\,\mathrm{s}$ per timestep on 27
ranks) a ten-year leg costs $1.6\,\mathrm{h}$ whatever follows it:

\begin{center}
\begin{tabular}{lrrrr}
\toprule
\textbf{Member} & \textbf{Forward leg} & \textbf{Adjoint} &
\textbf{Per member} & \textbf{300 members}\\
\midrule
1-year & $1.6$ h & $1.8$ h & $3.4$ node-h  & $\sim\!2.7\times10^{4}$ core-h\\
2-year & $1.6$ h & $3.6$ h & $5.2$ node-h  & $\sim\!4.2\times10^{4}$ core-h\\
5-year & $1.6$ h & $8.9$ h & $10.5$ node-h & $\sim\!8.5\times10^{4}$ core-h\\
\bottomrule
\end{tabular}
\end{center}

Adjoint times assume the gradient check is off. At one year the leg is $48\,\%$
of the member, so the saving from (2) runs out well short of what the adjoint
numbers alone suggest: a 300-member ensemble of one-year members lands near
$2.7\times10^{4}$ core-hours, against the $1.4\times10^{4}$ an adjoint-only count
gives. Members that vary only the background state are the exception --- they
start from an existing spin-up pickup, need no leg, and cost the adjoint alone,
which is an argument for weighting the ensemble towards the state axis wherever
the science allows.

These are still extrapolations from two measured rates and want replacing with a
real number. The first job in Phase~1 is to benchmark one one-year member with
\code{useGrdchk=.FALSE.} and report both halves separately.

\paragraph{Storage.}
A five-year run writes $8.3\,\mathrm{GB}$, mostly the $366$ dumps of eleven
fields. For training we would subsample the reverse sweep to about $12$ lead
times per member and keep \code{ADJdiffkr} plus the four auxiliary 3-D fields
($1.45\,\mathrm{MB}$ each in \code{float32}): five fields at twelve lead times
is $87\,\mathrm{MB}$ of 3-D labels per member, and the five 2-D forcing fields
add under $3\,\mathrm{MB}$ more. Across $300$ members that is
$\sim\!27\,\mathrm{GB}$, plus the forward state. Convert to Zarr at ingest with \code{xmitgcm} rather than
training off MDS directly.

\paragraph{One precision note.} \code{adxx\_*} files are \code{float64} and the
\code{ADJ*} diagnostic dumps are \code{float32}. Read the \code{.meta} instead
of assuming: guessing wrong reshapes the array into something that looks
plausible and is not.

%-----------------------------------------------------------------------
\section{Validation}
\label{sec:val}
%-----------------------------------------------------------------------

\paragraph{Metrics.} Volume-weighted pattern cosine similarity as the headline
number; $R^{2}$ in $\operatorname{asinh}$ space; log-amplitude error; relative
error of the scalar directional derivative
$\langle\hat{g},\delta\boldsymbol\theta\rangle$ over random and structured
perturbation directions; error in the leading POD coefficients; and an error
spectrum by horizontal wavenumber, which is how we would catch the
oversmoothing that regression losses on geophysical fields tend to produce.

\paragraph{Splits.} Five held-out axes, reported separately: (i) unseen
parameter draws inside the box; (ii) parameter-box corners, which is
extrapolation; (iii) unseen background states; (iv) an unseen cost section;
(v) unseen lead times.

\paragraph{Ground truth beyond the adjoint.}
The surrogate is trained on adjoint output, so on its own it can only inherit
whatever the adjoint gets wrong. Two independent checks:

\begin{enumerate}
  \item \textbf{Repair the gradient check first.} The built-in check has never
        passed, and the cause is where it is rather than the adjoint --- see
        Section~\ref{sec:validating} and Appendix~\ref{app:grdchk}. Until it is
        moved to $(2,127,26)$ with $\epsilon\approx10^{-3}$, the adjoint this
        surrogate would emulate is not verified by anything in the repository,
        and that has to be fixed before the rest of this is worth doing. Part~I's
        finite-difference comparison (Section~\ref{sec:fd}) is the stronger
        check of the same property.
  \item \textbf{Directional derivatives in $\kv$ directly.} Perturb the $\kv$
        field along a few smooth spatial directions, run paired forward models,
        and compare $\Delta J$ against $\langle g,\delta\kv\rangle$ from the
        adjoint and from the surrogate. This checks the quantity we are
        actually predicting rather than a state control.
\end{enumerate}

\paragraph{Downstream test.} The real test of a gradient surrogate is whether
it can be used for anything. Run a short calibration of $\kv$ against a target
transport twice, once with true adjoint gradients and once with surrogate
gradients, and compare the descent trajectories and the recovered fields. A
surrogate good for ten useful steps before it needs a true-adjoint correction
would already be a substantial win.

%-----------------------------------------------------------------------
\section{Risks and open questions}
\label{sec:risks}
%-----------------------------------------------------------------------

\begin{enumerate}
  \item \textbf{How far back the labels stay meaningful is not yet known.}
        The usual worry is that adjoint sensitivities in a chaotic flow grow
        exponentially backwards in time, and that \code{adjViscBoost} exists
        because long adjoints diverge. What I measured is more encouraging than
        that, and the $656\times$ growth quoted in Section~\ref{sec:adjoutput} should not be
        read
        as instability: that is \code{ADJdiffkr} accumulating, which is
        integration rather than amplification. The fields that genuinely
        propagate do the opposite. Over the same five years \code{ADJtheta} RMS
        \emph{decays} by about a factor of $5$ and \code{ADJtaux} by about
        $40$. At these viscosity settings the DINO state adjoint looks
        dissipation-dominated rather than explosive, which is good news for
        trainability.

        Two caveats. This is one run at one viscosity setting, and the
        \code{adjViscBoost} members exist because some configurations do
        misbehave, so the measurement should be repeated across the viscosity
        range before anyone leans on it. And decay has its own limit: once the
        signal has decayed into the noise floor the label is uninformative even
        though it is finite. Either way, $\tau_{\max}$ has to be established
        with directional-derivative checks at a sequence of lead times, before
        the ensemble is generated, because it sets the cost.
  \item \textbf{Two labelling functions.} Mixing boosted and unboosted adjoint
        sweeps without conditioning on \code{viscFacInAd} would corrupt the
        training set (Section~\ref{sec:inputs}).
  \item \textbf{Extrapolation.} We should claim nothing about parameter values
        outside the sampled box. Report the extrapolation split honestly and
        expect it to look bad.
  \item \textbf{Section indices are compiled in.} The cost-section axis needs
        either one rebuild per section, which is workable since builds are
        cheap next to runs, or better, moving
        \code{isecbeg/isecend/jsec/kmaxdepth} out of Fortran \code{parameter}
        statements and into \code{data.cost}. That is a single-file change that
        opens up a whole ensemble axis, so it is worth doing early.
  \item \textbf{GM/Redi is off}, so $\kappa_{\mathrm{GM}}$ and
        $\kappa_{\mathrm{Redi}}$ sensitivities are not available. Turning them
        on changes the physics and would need a new spin-up, which makes it a
        separate project and probably a follow-up paper rather than part of
        this one.
  \item \textbf{Idealised configuration.} DINO is one idealised basin. Whether
        a surrogate trained here carries over to a realistic configuration is
        an open question that we should not assume the answer to. What DINO
        buys us is an ensemble we can actually afford.
  \item \textbf{No Python environment is specified in the repository.} There is
        no environment file and no dependency manifest; the analysis notebooks
        assume \code{numpy}, \code{xarray}, \code{xmitgcm} and
        \code{matplotlib}. A pinned environment (PyTorch, Zarr, \code{xmitgcm})
        belongs in Phase 0 rather than being sorted out later.
\end{enumerate}

%-----------------------------------------------------------------------
\section{Roadmap}
%-----------------------------------------------------------------------

\begin{center}
\begin{tabularx}{\textwidth}{@{}Y{0.6}Y{0.55}Y{1.85}@{}}
\toprule
\textbf{Phase} & \textbf{Effort} & \textbf{Deliverable}\\
\midrule
0. Data audit
  & 2--3 weeks
  & Zarr ingest pipeline from MDS via \code{xmitgcm}; pinned Python
    environment; POD--ridge baseline fitted on the runs we already have;
    characterisation of the sensitivity field's statistics across lead time.\\
1. Configuration work
    and the scoping ensemble
  & 3--4 weeks
    $+$ 2--3 weeks
  & Repaired gradient check at $(2,127,26)$; \code{useGrdchk} off for ensemble
    members and the saving measured; cost-section indices moved into
    \code{data.cost}; a benchmarked one-year adjoint member with a real cost
    number; the trustworthy lead-time horizon $\tau_{\max}$ established.
    \textbf{Part~I runs here}: the seven-member $\kv$ ensemble, which answers
    whether mixing has to be an input at all and so gates Phase~2.\\
2. Ensemble generation
  & 4--8 weeks (mostly queue time)
  & $200$--$400$ member Sobol ensemble over parameters $\times$ background
    states $\times$ cost sections, ingested and split into
    train/validation/test along all five held-out axes.\\
3. Surrogate (F1)
  & 6--8 weeks
  & Trained conditioned U-Net; full metric table against the POD baseline on
    every split; error spectra; the directional-derivative validation.\\
4. Extensions
  & open
  & Sobolev-trained $J_\phi$ (F3); reverse-time propagator; $\Ah$ as a second
    field-valued target; the calibration demonstration.\\
\bottomrule
\end{tabularx}
\end{center}

\paragraph{The decision point.}
Phases 0 and 1 --- which include Part~I --- are worth doing whether or not the
surrogate ever gets built:
they repair the gradient check, cut the cost of every future adjoint run, and
make the cost functional configurable. The commitment comes at Phase~2, and
what should decide it is the benchmarked per-member cost from Phase~1 together
with how well the POD baseline does in Phase~0. If a reduced-basis model
already reproduces the parametric variation, the right thing to do is write
that up and skip the network.

